\section{Conclusion}

\label{sec:conclusion}

LP-Anchor provides a minimal, high-assurance mechanism for anchoring
CRDT state to a public blockchain. The design requires 72 bytes per anchor,
25,400 gas on C-Chain (near-zero on Q-Chain), and delivers sub-200ms
verification latency. The height monotonicity invariant and Quasar finality
guarantee that anchored state cannot be retroactively altered without
compromising the chain's validator set across three independent
cryptographic assumptions.

The precompile is general: any application that computes a periodic state
hash can use it. Base CRDT anchoring and \texttt{liquidity/state}
reconciliation anchoring are the first two consumers.

\bibliographystyle{plain}
\begin{thebibliography}{9}

\bibitem{quasar}
Z.~Kelling, ``LP-020: Quasar --- Post-Quantum BFT Consensus,''
\textit{Lux Research}, April 2026.

\bibitem{crdt}
M.~Shapiro et al., ``Conflict-free Replicated Data Types,''
\textit{Proc. SSS}, 2011.

\bibitem{ceramic}
Ceramic Network, ``Ceramic Protocol Specification v2,''
\url{https://ceramic.network}, 2024.

\bibitem{keri}
S.~Smith, ``Key Event Receipt Infrastructure (KERI),''
\textit{IETF Draft}, 2022.

\bibitem{eip2929}
V.~Buterin, M.~Swende, ``EIP-2929: Gas cost increases for state access opcodes,''
\textit{Ethereum Improvement Proposals}, 2021.

\end{thebibliography}

